\chapter*{Introduction}\label{ch:introduction}
\addcontentsline{toc}{chapter}{Introduction}

This report asks a deceptively simple question: \emph{Why do we sleep, and what are dreams for?} 

Sleep is a nearly universal behavior among the Earth's animals, yet its function remains enigmatic. Dreams, the vivid experiences that occur during sleep, have been variously interpreted as meaningless byproducts or as having specific psychological functions.

We argue that sleep is a computational necessity for any intelligence capable of learning: an offline, staged \emph{learning algorithm} that adaptive nervous systems require. Dreams are neither useless noise nor the purpose of sleep; they are its most vivid signature: the subjective experience of memory consolidation, emotional processing, and cognitive generalization while it happens. 

Part~\ref{part:culture} surveys what human cultures intuited about dreams; Part~\ref{part:science} outlines what is established about sleep`s architecture and function; Part~\ref{part:obh} integrates these into a unified account, the \emph{Overfitted-Brain Hypothesis}.
